\documentclass[12pt]{article}
\usepackage[utf8]{inputenc}
\usepackage[T1]{fontenc}
\usepackage[margin=1in]{geometry}
\usepackage{setspace}
\usepackage{xcolor}
\usepackage[colorlinks=true, urlcolor=blue!60!black, linkcolor=blue!60!black]{hyperref}
\setlength{\parindent}{2em}
\setlength{\parskip}{1.25em}
\renewcommand{\baselinestretch}{1.0}
\begin{document}
\begin{center}
{\LARGE \textit{Kennedy v. Bremerton School District}, 597 U.S. 507 (2022)\par}
\vspace{0.8em}
{\large Case Brief --- Four-Page Edition\par}
\end{center}
\vspace{1em}

\section{Memory Jogger}

A public high school football coach's brief, quiet postgame prayer at midfield is protected by the Free Exercise and Free Speech Clauses, and the Establishment Clause, interpreted by historical practices and understandings rather than the abandoned \textit{Lemon} and endorsement tests, neither requires nor permits the school district to suppress it (\href{https://www.law.cornell.edu/supremecourt/text/21-418}{597 U.S. 507}).

\section{Detailed Case Facts}

Joseph Kennedy, a Christian and Marine Corps veteran, coached football at Bremerton High School in Washington beginning in 2008. After each game he knelt at the 50-yard line for a quiet prayer of about 30 seconds. Over time players joined him (``This is a free country. You can do what you want,'' he told them), and the practice grew into motivational talks with religious references; he also took part in an existing tradition of locker-room prayer.

In September 2015 the District directed that any on-duty religious expression be ``nondemonstrative.'' Kennedy stopped the locker-room prayers and religious talks, but his lawyers then announced that he would resume a short, private midfield prayer, and he publicized his plans in the media. The District barred ``any overt actions'' that a reasonable observer would see as endorsing prayer. After three October games, while players sang the fight song, Kennedy knelt and prayed; no Bremerton players joined him, though spectators rushed the field after one game. The District placed him on paid leave, conceding it had no evidence that students were directly coerced, and he did not return.

\section{Procedural History}

Kennedy sued under the Free Speech and Free Exercise Clauses and Title VII. The district court denied a preliminary injunction and the Ninth Circuit affirmed; the Supreme Court denied certiorari in 2019 over a statement by Justice Alito. On remand, the district court granted summary judgment to the District, and the Ninth Circuit affirmed, holding the prayers government speech and, alternatively, that avoiding an endorsement violation justified the suspension; rehearing en banc was denied over 11 dissents. The Supreme Court reversed on June 27, 2022.

\section{Judicial Votes}

Reversed, \textbf{6--3}.
\begin{itemize}
\item \textbf{Majority:} Justice Gorsuch, joined by Chief Justice Roberts and Justices Thomas, Alito, and Barrett, and by Justice Kavanaugh except Part III--B.
\item \textbf{Concurring:} Justice Thomas; Justice Alito.
\item \textbf{Dissenting:} Justice Sotomayor, joined by Justices Breyer and Kagan.
\end{itemize}

\section{Holding}

The Free Exercise and Free Speech Clauses protect a public employee's brief, quiet, personal religious observance during a postgame lull when employees may attend to personal matters. The Establishment Clause must be interpreted by historical practices and understandings; \textit{Lemon v.\ Kurtzman} and its endorsement offshoot are abandoned. Absent coercion, the District could not suppress Kennedy's prayers. Whether \textit{Smith} should be reconsidered and what test governs employees' free-exercise claims were left open.

\section{Analysis of Opinions}

\textbf{Majority (Gorsuch).} The Free Exercise and Free Speech Clauses work in tandem. The District's policies were not neutral, since suppressing religious conduct was their object, nor generally applicable, since other coaches could make phone calls or greet friends during the same period. Under \textit{Pickering}--\textit{Garcetti} and \textit{Lane v.\ Franks}, the prayers did not owe their existence to Kennedy's duties: he was not instructing players or delivering a commissioned message. Treating all on-duty speech as government speech would let schools ban a Muslim teacher's headscarf.

The Clauses are complementary, not warring. The Court had long abandoned \textit{Lemon}, an ``ambitious, abstract, and ahistorical'' test that worked as a modified heckler's veto. Coercion is a hallmark of establishment, but none was shown: no students joined, the District conceded there was no evidence of direct coercion, and visibility or offense is not coercion. \textit{Lee v.\ Weisman} and \textit{Santa Fe v.\ Doe} involved captive audiences or school-sponsored messages.

\textbf{Thomas and Alito (concurring).} Thomas flagged whether employees' free-exercise rights differ from the public's and what scrutiny applies. Alito read the opinion as leaving open the free-speech standard for private expression during a lull.

\textbf{Sotomayor (dissenting).} The majority ``misconstrues the facts'' by isolating three prayers from seven years of audible midfield prayers, a media campaign, and disruption on the field. Under \textit{Engel}, \textit{Schempp}, \textit{Lee}, and \textit{Santa Fe}, schools must remain neutral because students face subtle pressure from coaches who control playing time. The new history-and-tradition test offers schools ``essentially no guidance for school administrators.''

\section{Examples: Future Applications}

\begin{enumerate}
\item \textbf{Same-Side (Protected): Lunchtime Prayer.} A math teacher silently prays at her desk during lunch while students sometimes pass the open door. It is private exercise on personal time without coercion.
\item \textbf{Same-Side (Protected): Religious Attire.} A coach wears a yarmulke or hijab on the sidelines where secular hats are allowed. A ban would be neither neutral nor generally applicable.
\item \textbf{Opposite-Side (Unprotected): Intercom Prayer.} A principal opens mandatory homeroom by reciting a prayer over the intercom. Official speech to a captive audience is barred by \textit{Engel}, \textit{Lee}, and \textit{Santa Fe}.
\item \textbf{Opposite-Side (Unprotected): Coach-Led Huddle.} A coach leads an audible prayer in a mandatory postgame huddle, crossing the coercion line \textit{Kennedy} preserved.
\item \textbf{Fence-Sitter: Prayer Before the Critique.} A track coach kneels silently at the finish line, knowing runners routinely gather to pray, then immediately delivers a mandatory critique. He invites no one, but the transition and recurring participation suggest indirect coercion.
\end{enumerate}

\section{Critique}

\textbf{Originalist and conservative views.} Supporters praise the burial of \textit{Lemon} and the harmonizing of the three Clauses. Stephanie Taub and Kayla Ann Toney argue that historical practices, informed by Michael McConnell's Founding-era hallmarks of establishment, anchor doctrine in original meaning. Some originalist commentators note, however, that the opinion is thin on Founding-era analysis and does not explain how to analogize eighteenth-century establishments to modern public schools.

\textbf{Progressive views.} Ira Lupu and Robert Tuttle call the decision ``a sledgehammer to the bedrock of nonestablishment,'' arguing that it distorted the record and hollowed out the School Prayer Cases. Reducing protection to proof of direct coercion ignores the power dynamics between coaches and adolescents and leaves minority-faith students exposed.

\section{Key Quotations}

\begin{quote}
``In place of \textit{Lemon} and the endorsement test, this Court has instructed that the Establishment Clause must be interpreted by `reference to historical practices and understandings.' ''
\end{quote}
\noindent --- Justice Gorsuch, Opinion of the Court, 597 U.S. 507.

\begin{quote}
``The Constitution and the best of our traditions counsel mutual respect and tolerance, not censorship and suppression, for religious and nonreligious views alike.''
\end{quote}
\noindent --- Justice Gorsuch, Opinion of the Court, 597 U.S. 507.

\begin{quote}
``And in no world may a government entity's concerns about phantom constitutional violations justify actual violations of an individual's First Amendment rights.''
\end{quote}
\noindent --- Justice Gorsuch, Opinion of the Court, 597 U.S. 507.

\end{document}
